\documentclass[10pt,a4paper]{amsart}
\usepackage{amsmath,amssymb,amsthm,graphicx}
% Retained minimal preamble for the Stacks Project chapter more-algebra.tex (More on Algebra).
% The chapter file pinned-original-source.tex carries only \input{preamble} on its first line, so no
% preamble text lives inside the chapter itself.  The Stacks Project uses one single shared preamble for
% every chapter; the snapshot stages it identically as sources/stacks-chow/preamble.tex and
% sources/stacks-etale-cohomology/preamble.tex.  The block below is copied mechanically, line for line,
% from sources/stacks-chow/preamble.tex lines 14-18 (the Xy-pic setup) and 154-217 (the native theorem
% declarations together with the mathematical macros).  Nothing was added, removed or redefined; the
% wrapper class and the amsmath/amssymb/amsthm/graphicx base packages are supplied by the assembler.
% Chapter: more-algebra.tex -- More on Algebra.
% For commutative diagrams we use Xy-pic
\usepackage[all]{xy}

% We use 2cell for 2-commutative diagrams.
\xyoption{2cell}

% Theorem environments.
%
\theoremstyle{plain}
\newtheorem{theorem}[subsection]{Theorem}
\newtheorem{proposition}[subsection]{Proposition}
\newtheorem{lemma}[subsection]{Lemma}

\theoremstyle{definition}
\newtheorem{definition}[subsection]{Definition}
\newtheorem{example}[subsection]{Example}
\newtheorem{exercise}[subsection]{Exercise}
\newtheorem{situation}[subsection]{Situation}

\theoremstyle{remark}
\newtheorem{remark}[subsection]{Remark}
\newtheorem{remarks}[subsection]{Remarks}

\numberwithin{equation}{subsection}

% Macros
%
\def\lim{\mathop{\mathrm{lim}}\nolimits}
\def\colim{\mathop{\mathrm{colim}}\nolimits}
\def\Spec{\mathop{\mathrm{Spec}}}
\def\Hom{\mathop{\mathrm{Hom}}\nolimits}
\def\Ext{\mathop{\mathrm{Ext}}\nolimits}
\def\SheafHom{\mathop{\mathcal{H}\!\mathit{om}}\nolimits}
\def\SheafExt{\mathop{\mathcal{E}\!\mathit{xt}}\nolimits}
\def\Sch{\mathit{Sch}}
\def\Mor{\mathop{\mathrm{Mor}}\nolimits}
\def\Ob{\mathop{\mathrm{Ob}}\nolimits}
\def\Sh{\mathop{\mathit{Sh}}\nolimits}
\def\NL{\mathop{N\!L}\nolimits}
\def\CH{\mathop{\mathrm{CH}}\nolimits}
\def\proetale{{pro\text{-}\acute{e}tale}}
\def\etale{{\acute{e}tale}}
\def\QCoh{\mathit{QCoh}}
\def\Ker{\mathop{\mathrm{Ker}}}
\def\Im{\mathop{\mathrm{Im}}}
\def\Coker{\mathop{\mathrm{Coker}}}
\def\Coim{\mathop{\mathrm{Coim}}}

% Boxtimes
%
\DeclareMathSymbol{\boxtimes}{\mathbin}{AMSa}{"02}

%
% Macros for moduli stacks/spaces
%
\def\QCohstack{\mathcal{QC}\!\mathit{oh}}
\def\Cohstack{\mathcal{C}\!\mathit{oh}}
\def\Spacesstack{\mathcal{S}\!\mathit{paces}}
\def\Quotfunctor{\mathrm{Quot}}
\def\Hilbfunctor{\mathrm{Hilb}}
\def\Curvesstack{\mathcal{C}\!\mathit{urves}}
\def\Polarizedstack{\mathcal{P}\!\mathit{olarized}}
\def\Complexesstack{\mathcal{C}\!\mathit{omplexes}}
% \Pic is the operator that assigns to X its picard group, usage \Pic(X)
% \Picardstack_{X/B} denotes the Picard stack of X over B
% \Picardfunctor_{X/B} denotes the Picard functor of X over B
\def\Pic{\mathop{\mathrm{Pic}}\nolimits}
\def\Picardstack{\mathcal{P}\!\mathit{ic}}
\def\Picardfunctor{\mathrm{Pic}}
\def\Deformationcategory{\mathcal{D}\!\mathit{ef}}

% --- appended 2026-10-12 (Chief): the project own definitions of the reference/slogan/history
% display environments, copied mechanically from the shared Stacks preamble lines 8-12
% (\usepackage{verbatim} plus the three \newenvironment lines, which the chapter preamble keeps
% verbatim from the project preamble). Nothing added, removed or redefined: these environments
% carry no mathematical content and the project itself typesets them as comments. Without these
% three lines a counted statement whose own text uses them fails with Environment slogan undefined.
\usepackage{verbatim}
\newenvironment{reference}{\comment}{\endcomment}
\newenvironment{slogan}{\comment}{\endcomment}
\newenvironment{history}{\comment}{\endcomment}
\usepackage{cleveref}
\crefname{originalproposition}{original proposition}{original propositions}
\Crefname{originalproposition}{Original proposition}{Original propositions}
\makeatletter
\newlabel{proposition-noetherian-naive-completion-is-completion}{{0922}{source}{}{sourcepointer-1}{}}
\newlabel{proposition-noetherian-naive-completion-is-completion@cref}{{[originalproposition][0][]0922}{[1][0][]source}{}{}{}}
\makeatother
\crefname{originallemma}{original lemma}{original lemmas}
\Crefname{originallemma}{Original lemma}{Original lemmas}
\makeatletter
\newlabel{lemma-derived-completion-pseudo-coherent}{{0A06}{source}{}{sourcepointer-2}{}}
\newlabel{lemma-derived-completion-pseudo-coherent@cref}{{[originallemma][0][]0A06}{[1][0][]source}{}{}{}}
\makeatother
\crefname{originallemma}{original lemma}{original lemmas}
\Crefname{originallemma}{Original lemma}{Original lemmas}
\makeatletter
\newlabel{lemma-break-long-exact-sequence-modules}{{0CQE}{source}{}{sourcepointer-3}{}}
\newlabel{lemma-break-long-exact-sequence-modules@cref}{{[originallemma][0][]0CQE}{[1][0][]source}{}{}{}}
\makeatother
\crefname{originallemma}{original lemma}{original lemmas}
\Crefname{originallemma}{Original lemma}{Original lemmas}
\makeatletter
\newlabel{lemma-compute-Rlim-modules}{{091D}{source}{}{sourcepointer-4}{}}
\newlabel{lemma-compute-Rlim-modules@cref}{{[originallemma][0][]091D}{[1][0][]source}{}{}{}}
\makeatother
\title{The limit of the cohomology of the quotients computes the cohomology of the derived completion}
\author{The Stacks Project Authors}
\date{stacks-project:more-algebra.tex:local-snapshot-20261012; GNU FDL 1.2 or later}
\begin{document}
\maketitle
\setcounter{subsection}{1}
\noindent\textit{Editorial note.} All delivered mathematical text is the author's own, copied from the pinned Stacks Project chapter file more-algebra.tex; no statement or proof step was rewritten, reordered or removed. The counted span is the author environment and its own proof environment. References inside the span are handed to the assembler as navigation stubs that carry the canonical tag of the cited tag only; whether each of them is load-bearing is recorded on the stub itself.\par\medskip
\noindent\textit{Allowed standard prerequisites (editorial).}\begin{itemize}
\item For a Noetherian ring the derived completion of a complex is the derived limit of the complexes tensored with the quotients by the powers of the ideal. invoked at proof lines 28538-28540. Identifies the object whose cohomology is being computed with the derived limit of the systems appearing in the statement.
\item The cohomology of the derived completion of a pseudo-coherent complex is the inverse limit of the cohomology of the quotients. invoked at proof lines 28541-28542. Computes the cohomology of the derived completion as the inverse limit of the cohomology of the powers, which is the left hand side of the statement.
\item A short exact sequence of inverse systems gives a short exact sequence relating the derived limit and the inverse limit, with the obstruction term on the left. invoked at proof lines 28543-28547. Produces the displayed short exact sequence whose middle term is the computed cohomology and whose right hand term is the limit in the statement.
\end{itemize}\medskip
\ifdefined\hypertarget\hypertarget{sourcepointer-1}{}\fi
\noindent\textit{Source reference.} Original proposition 0922: sources/stacks-more-algebra/pinned-original-source.tex, source lines 27219--27219. Invoked at proof lines 28538-28540 to identify the derived completion; the computation of its cohomology and the exact sequence argument that follow are inside the counted span.\par\medskip
\ifdefined\hypertarget\hypertarget{sourcepointer-2}{}\fi
\noindent\textit{Source reference.} Original lemma 0A06: sources/stacks-more-algebra/pinned-original-source.tex, source lines 27266--27266. Used at proof lines 28541-28542 to identify the cohomology of the derived completion with the inverse limit of the cohomology of the quotients; the rest of the argument is inside the counted span.\par\medskip
\ifdefined\hypertarget\hypertarget{sourcepointer-3}{}\fi
\noindent\textit{Source reference.} Original lemma 0CQE: sources/stacks-more-algebra/pinned-original-source.tex, source lines 23653--23653. Applied at proof lines 28543-28547 to obtain the displayed exact sequence; the vanishing argument using Mittag-Leffler that finishes the proof is inside the counted span.\par\medskip
\ifdefined\hypertarget\hypertarget{sourcepointer-4}{}\fi
\noindent\textit{Source reference.} Original lemma 091D: sources/stacks-more-algebra/pinned-original-source.tex, source lines 23574--23574. Used at proof lines 28548-28549 to see that the left hand term of the exact sequence vanishes; the conclusion drawn from that vanishing is inside the counted span.\par\medskip
\begin{lemma}[Koll\'ar-Kov\'acs]
\label{lemma-kollar-kovacs}
\begin{reference}
Email from Kovacs of 23/02/2018.
\end{reference}
Let $I$ be an ideal of a Noetherian ring $A$. Let $K \in D(A)$.
Set $K_n = K \otimes_A^\mathbf{L} A/I^n$. Assume for all
$i \in \mathbf{Z}$ we have
\begin{enumerate}
\item $H^i(K)$ is a finite $A$-module, and
\item the system $H^i(K_n)$ satisfies Mittag-Leffler.
\end{enumerate}
Then $\lim H^i(K)/I^nH^i(K)$ is equal to $\lim H^i(K_n)$ for all
$i \in \mathbf{Z}$.
\end{lemma}

\begin{proof}
Recall that $K^\wedge = R\lim K_n$ is the derived completion of $K$, see
Proposition \ref{proposition-noetherian-naive-completion-is-completion}.
By Lemma \ref{lemma-derived-completion-pseudo-coherent} we have
$H^i(K^\wedge) = \lim H^i(K)/I^nH^i(K)$. By
Lemma \ref{lemma-break-long-exact-sequence-modules}
we get short exact sequences
$$
0 \to R^1\lim H^{i - 1}(K_n) \to H^i(K^\wedge) \to \lim H^i(K_n) \to 0
$$
The Mittag-Leffler condition guarantees that the left terms are zero
(Lemma \ref{lemma-compute-Rlim-modules}) and we conclude the lemma is true.
\end{proof}

\end{document}
